\addtocontents{toc}{\protect\Needspace{3\baselineskip}}
\originalchapter{体场、尺度维数与关联函数}\label{ch:correlator}
\originalsection{质量与维数}
体标量作用量 $S=-\frac12\int\sqrt{-g}[(\nabla\phi)^2+m^2\phi^2]$ 给 $(\Box-m^2)\phi=0$. Poincaré 坐标在 $z\to0$ 忽略有限边界动量项, 得 $z^2\phi''+(1-d)z\phi'-m^2R^2\phi=0$. 代 $\phi\sim z^s$ 得 $s(s-d)=m^2R^2$, 根 $\Delta_\pm=d/2\pm\nu$, $\nu=\sqrt{d^2/4+m^2R^2}$. 标准量子化使用 $\Delta=\Delta_+$, 边界展开为 $\phi=z^{d-\Delta}\phi_0+z^\Delta A+\cdots$.

实根要求 Breitenlohner--Freedman 下界 $m^2R^2\ge-d^2/4$. AdS 可允许一定负质量平方, 不能按 Minkowski tachyon 判据直接断言不稳定. 在 $0<\nu<1$ 的特定窗口可有替代量子化, 但需要重新选择源、期望值和边界作用量, 本卷主要采用标准量子化. 根重合或整数差时可能出现对数项, 不能盲目把两个纯幂当成完整展开.

\originalsection{欧氏边界到体核}
令 $R=1$, 欧氏度规 $\dd s^2=(\dd z^2+\dd x^2)/z^2$. 对标准维数 $\Delta>d/2$, 正则边界到体核为
\begin{equation}\label{eq:K}
K(z,x;x')=C_\Delta\left(\frac z{z^2+|x-x'|^2}\right)^\Delta,
\quad C_\Delta=\frac{\Gamma(\Delta)}{\pi^{d/2}\Gamma(\Delta-d/2)}.
\end{equation}
直接微分可验证场方程; 归一化来自 $\int\dd^dy(1+y^2)^{-\Delta}=\pi^{d/2}\Gamma(\Delta-d/2)/\Gamma(\Delta)$, 保证 $z^{\Delta-d}\phi(z,x)\to\phi_0(x)$ 的分布极限. 不能对固定 $x\ne x'$ 的点态极限替代这个 delta 核极限.

积分分部使壳上二次作用量变为截止面边界项. 去掉局域发散项后, 非局域两点核为
$\langle\mathcal O(x)\mathcal O(0)\rangle=(2\Delta-d)C_\Delta/|x|^{2\Delta}$, 此式取体动能规范化为 $1/2$. 若整体作用量有 $1/\kappa^2$, 关联函数还带该系数. 接触项依反项方案改变, 分离点幂律不变.

\originalsection{高点函数}
加入体相互作用 $g_3\phi^3/3!$, 三点函数在树级为 $-g_3\int\dd^{d+1}X\sqrt g\,K(X;x_1)K(X;x_2)K(X;x_3)$, 符号随源与欧氏作用量约定保持一致. 它有局域体顶点, 但在边界是非局域积分核. 树级交换图还需体到体传播子; 环图、混合算符和全套全息重整化不在本卷已完整计算的范围. 本卷 PS5.2 的题解证明核的基本极限关系, 不把所有高点积分都标成已求闭式.

\originalsection{两点核的壳上计算}
边界到体核的归一化可逐步完成. 令 $r=|x-x'|$, 对 $K=C(z/(z^2+r^2))^\Delta$ 的对数求导给 $\partial_z\ln K=\Delta(1/z-2z/(z^2+r^2))$, $\partial_i\ln K=-2\Delta(x_i-x'_i)/(z^2+r^2)$. 把 $\partial^2K/K=(\partial\ln K)^2+\partial^2\ln K$ 代进 AdS Laplacian, 剩余常数为 $\Delta(\Delta-d)$, 正好等于 $m^2R^2$. 这实际验证了场方程, 而非只凭尺度形状猜核.

归一化积分取 $x-x'=zy$, 源解 $\phi(z,x)=\int K\phi_0\dd^dx'$ 的领先项为 $Cz^{d-\Delta}\phi_0(x)\int\dd^dy(1+y^2)^{-\Delta}$. 径向积分再令 $u=r^2$, 得 $\int=\pi^{d/2}\Gamma(\Delta-d/2)/\Gamma(\Delta)$, 确定 $C=C_\Delta$. 对分离点, 响应系数 $A(x)=C_\Delta\int\phi_0(x')\dd^dx'/|x-x'|^{2\Delta}$; 对重合点须作分布延拓并减局域项.

以欧氏正动能作用量积分分部, 壳上边界项是 $S_E=-R^{d-1}\int\epsilon^{1-d}\phi\partial_z\phi\dd^dx/2$. 设 $\phi=\epsilon^{d-\Delta}\phi_0+\epsilon^\Delta A+\cdots$. 领先发散项被局域反项 $(d-\Delta)\int\sqrt\gamma\phi^2/(2R)$ 抵消. 非局域交叉项则合并为
$S_{E,\rm ren}^{\rm nonlocal}=-(\Delta-d/2)R^{d-1}\int\phi_0A\dd^dx$.
当存在额外导数发散时还需相应局域反项, 它们对分离点只改变接触项. 由 $W=-S_E$ 对两个源求导, 响应核对称给另一个因子 2, 得 $(2\Delta-d)R^{d-1}C_\Delta/|x-x'|^{2\Delta}$. 正文结果取 $R=1$, 并已固定体场动能归一化.
